% !TEX root = ../main.tex
\chapter{Kết luận}
\label{Chapter5}

Chương này tổng kết mức độ hoàn thành mục tiêu của ConferenceSpace bằng cách đối chiếu với vấn đề nêu ở Chương~1, kết quả khảo sát và quy ước trách nhiệm ở Chương~2, thiết kế ở Chương~3 cùng bằng chứng thực nghiệm ở Chương~4.

Trong điều kiện thử nghiệm, ConferenceSpace đã triển khai mô hình phân tách trách nhiệm thành ba lớp và thu thập bằng chứng cho từng tác vụ hỗ trợ được đánh giá. Tuy nhiên, báo cáo không kết luận rằng mô hình này duy trì hiệu quả nội dung phản biện và bảo đảm mọi quyết định vẫn thuộc về con người. Reviewer Initial Analysis tạo nội dung có thể ảnh hưởng đến nhận định, còn Review Quality Auditor cho phép một phát hiện do AI tạo ở mức \textit{blocking} ngăn thao tác gửi trong khi Phản biện viên chưa có quyền ghi đè. Báo cáo cũng không khẳng định mô hình ba lớp là kiến trúc tối ưu hoặc ConferenceSpace ưu việt hơn các nền tảng lấy AI làm trọng tâm đã khảo sát. Đóng góp được xác nhận trong phạm vi đề tài là việc triển khai mô hình trách nhiệm, đồng thời xác định bằng chứng và khoảng trống của từng tác vụ.

\section{Kết quả đạt được}

\subsection{Nền tảng nghiệp vụ trong phạm vi đề tài}

Đề tài đã xây dựng các thành phần, giao diện và API phục vụ vòng đời bình duyệt theo các vai trò Tác giả, Phản biện viên và Chủ tọa/Đồng chủ tọa. Phạm vi gồm cấu hình hội nghị và chuyên đề, nộp và sửa bài, phân công, kiểm tra xung đột lợi ích, phản biện, phản hồi của Tác giả, thảo luận, ra quyết định và gửi thông báo. Đối với bản hoàn thiện sau chấp nhận, hệ thống mới hỗ trợ tải lên và truy xuất tệp; hệ thống chưa hỗ trợ Chủ tọa phê duyệt, kiểm soát hạn chót hoặc yêu cầu Tác giả nộp lại. Chương~4 chỉ xác nhận hiệu năng của một số endpoint; thử nghiệm chưa kiểm tra đầu cuối tính đúng đắn của toàn bộ chuỗi trạng thái và cơ chế phân quyền.

Ba kịch bản k6 ghi nhận tỷ lệ lỗi yêu cầu bằng 0\% và độ trễ p95 dưới 120~ms trên các endpoint đã chọn. Kết quả này xác nhận các endpoint đáp ứng ngưỡng hiệu năng trong cấu hình thử nghiệm, nhưng không thay thế phép kiểm thử chức năng cho toàn bộ vòng đời và không chứng minh hệ thống đã sẵn sàng phục vụ một hội nghị thực tế. Các chức năng quản lý sự kiện ngoài trọng tâm bình duyệt, mức độ hoàn thiện tương đương sản phẩm thương mại và khả năng vận hành dài hạn vẫn nằm ngoài phạm vi hoặc chưa được triển khai đầy đủ.

\subsection{Cách phân công trách nhiệm giữa ba lớp}

Đề tài đã triển khai mô hình phân tách trách nhiệm thành ba lớp logic. Lớp nghiệp vụ cốt lõi quản lý trạng thái, quyền hạn và quy tắc của vòng đời hội nghị. Lớp thuật toán xử lý quy tắc nộp bài, xếp hạng bằng chỉ số Jaccard và một phần hoạt động kiểm tra xung đột lợi ích. Phép xếp hạng Jaccard cho kết quả xác định, nhưng quy trình phân công có lượt dự phòng nới lỏng ngưỡng nên không bảo đảm mức độ phù hợp chuyên môn. Lớp AI tạo bản nháp, cảnh báo, nội dung phân tích hoặc bản tổng hợp để người có thẩm quyền xem xét. Sự phân tách này xác định ranh giới trách nhiệm về mặt logic; thử nghiệm chưa chứng minh mọi ranh giới đều được tuân thủ hiệu quả trong điều kiện sử dụng thực tế.

Trên tập dữ liệu tổng hợp mô phỏng lược đồ Semantic Scholar, Jaccard truy hồi hồ sơ của tác giả gốc tốt hơn phương pháp ngẫu nhiên trong phép thử theo chủ đề, với MRR đạt 0,392, Hit@5 đạt 55\% và Hit@10 đạt 65\%. Phép thử này chưa đo mức độ phù hợp chuyên môn của một Phản biện viên độc lập. Trong quy trình phân công, tỷ lệ bài được phân đủ hai Phản biện viên chỉ đạt 65,9\%; 23,3\% số bài phải dùng lượt dự phòng, trong đó cả ngưỡng điểm và giới hạn tải tối đa đều bị loại bỏ. Vì vậy, kết quả chỉ mô tả khả năng truy hồi theo chủ đề và hành vi nội tại của quy trình phân công; bằng chứng hiện có chưa đủ để hệ thống tự động xác nhận phân công hoặc để so sánh chất lượng với nền tảng khác.

Cơ chế phát hiện xung đột lợi ích kết hợp kiểm tra trường hợp tự phân công, thông tin khai báo và dữ liệu đồng tác giả. Microbenchmark chỉ đánh giá hai cơ chế phát hiện trường hợp tự phân công và xung đột đã khai báo; báo cáo chưa trình bày kết quả microbenchmark đối với lớp quan hệ đồng tác giả trên Neo4j. Tập dữ liệu phân công không có trường hợp tự phân công, trong khi tất cả phương pháp cơ sở đều cho cùng kết quả và chưa có tập các cặp xung đột thuộc nhiều loại do chuyên gia gán nhãn. Do đó, báo cáo chưa cung cấp bằng chứng về chất lượng của cơ chế phát hiện xung đột lợi ích nhiều lớp.

\subsection{Chuỗi bằng chứng cho các chức năng AI hỗ trợ}

Đề tài không dựa vào số lượng chức năng AI để kết luận về chất lượng hệ thống. Mỗi luồng xử lý có nhiệm vụ, đầu ra, người chịu trách nhiệm xác nhận và phương pháp đánh giá riêng.

Submission Autofill có bằng chứng trực tiếp rõ nhất đối với siêu dữ liệu: F1 theo token tiêu đề đạt 98,20\%, F1 của từ khóa đạt 92,77\% và tỷ lệ hoàn tất các trường bắt buộc đạt 86,93\%. Giá trị thấp nhất bằng 0 ở một số trường và F1 của thông tin tác giả thấp hơn cho thấy Tác giả vẫn phải kiểm tra bản nháp. Submission Gating xử lý đúng 8/8 trường hợp kiểm tra quy tắc trong bộ thử hiện tại và giữ 24 trường hợp kiểm tra bằng AI ở mức cảnh báo; chất lượng nội dung của các cảnh báo này chưa được chuyên gia gán nhãn.

Theo bộ chấm tự động mang tính thăm dò, Reviewer Initial Analysis đạt tỷ lệ trích dẫn bám nguồn 96,22\%, nhưng độ trung thực của các điểm cần chú ý chỉ đạt 69,86\%. Luồng này còn tạo điểm mạnh, điểm yếu và gợi ý có thể ảnh hưởng đến nội dung phản biện; thử nghiệm chưa đánh giá mức độ Phản biện viên phụ thuộc vào nội dung gợi ý hoặc mức độ phù hợp của chức năng với chính sách hội nghị.

Review Quality Auditor chỉ đạt 46,99\% phát hiện vừa bám nguồn vừa hợp lệ. Trong khi đó, một phát hiện do AI tạo ở mức \textit{blocking} có thể ngăn Phản biện viên gửi bản phản biện, nhưng Phản biện viên chưa có quyền ghi đè. Vì vậy, luồng này chưa đáp ứng yêu cầu bảo đảm quyền kiểm soát của con người. Các phát hiện do AI tạo cần được chuyển thành cảnh báo; chỉ các quy tắc xác định mới được phép chặn thao tác gửi.

Chair Decision Copilot đạt 87,34\% về độ trung thực của cơ sở bằng chứng và 87,11\% về độ trung thực của bản tổng hợp các điểm bất đồng theo phép chấm TCA. Kiểm thử chấp nhận của người dùng (User Acceptance Testing/UAT) dành cho Chủ tọa cung cấp tín hiệu cảm nhận ban đầu về tính minh bạch và khả năng chất vấn các công cụ AI nói chung, nhưng không đánh giá riêng chức năng này hoặc đo ảnh hưởng của nó đến chất lượng quyết định. Các phép chấm TCA và NLI chưa được hiệu chuẩn bằng nhãn chuyên gia, nên báo cáo chỉ dùng chúng làm chỉ số gián tiếp để lựa chọn trường hợp cần kiểm tra.

Chatbot Agent hoàn tất 40 hội thoại và không làm lộ dữ liệu trong các kịch bản kiểm tra quyền truy cập đã thực hiện. Tuy nhiên, một biến thể của yêu cầu nghiên cứu ngoài nền tảng đã tạo báo cáo vượt phạm vi; chỉ 37/40 hội thoại được đánh giá là đạt hoặc đạt một phần, còn 31/128 lượt gọi công cụ thất bại. Kết quả cho thấy kiểm soát quyền truy cập và kiểm soát phạm vi nội dung là hai thuộc tính khác nhau. Các kịch bản đã thực hiện không tương đương một cuộc kiểm toán bảo mật và cũng không cung cấp bằng chứng về độ ổn định trong vận hành dài hạn.

\subsection{Mức độ hoàn thành mục tiêu}

ConferenceSpace đã xây dựng một nền tảng thực nghiệm áp dụng mô hình phân tách trách nhiệm thành ba lớp và tạo chuỗi bằng chứng ban đầu cho từng tác vụ được đánh giá. Độ vững của bằng chứng không đồng đều: hiệu năng của các endpoint đã chọn, độ chính xác khi trích xuất siêu dữ liệu của Submission Autofill và bộ trường hợp kiểm tra quy tắc có phép đo trực tiếp; kết quả truy hồi theo chủ đề và phép chấm TCA chỉ là chỉ số gián tiếp; còn UAT chỉ phản ánh cảm nhận của những người tham gia khảo sát. Thử nghiệm chưa đo đầy đủ tính đúng đắn đầu cuối, chất lượng của cơ chế phát hiện xung đột lợi ích nhiều lớp, hiệu lực kiểm soát của con người, hoạt động quản trị dữ liệu của nhà cung cấp và khả năng chịu lỗi.

Do đó, kết quả chỉ hỗ trợ kết luận rằng ConferenceSpace đã triển khai các ràng buộc không cho AI tự điền hoặc gửi biểu mẫu phản biện, chấm điểm hay đưa ra quyết định chấp nhận hoặc từ chối bài. Thử nghiệm chưa đo hiệu lực thực tế của việc người dùng kiểm tra, bác bỏ hoặc ghi đè đầu ra AI. Reviewer Initial Analysis vẫn tạo nội dung tương tự một phần của bản phản biện, còn Review Quality Auditor vượt quá vai trò tư vấn khi một phát hiện do AI tạo có thể chặn thao tác gửi. Đề tài chưa chứng minh ConferenceSpace ưu việt hơn nền tảng khác, chưa xác lập quan hệ nhân quả giữa mô hình ba lớp và kết quả quan sát được, đồng thời chưa chứng minh hệ thống sẵn sàng triển khai thực tế.

\section{Các hạn chế}

\subsection{Giới hạn của so sánh và phạm vi sản phẩm}

Phần so sánh ở Chương~2 dựa trên tài liệu sản phẩm, tài liệu hỗ trợ, bài viết kỹ thuật và chính sách công khai. Nhóm không so sánh trực tiếp ConferenceSpace với PeerSubmit, Morressier hoặc nền tảng thương mại khác trên cùng dữ liệu và bằng cùng phương pháp. Vì vậy, báo cáo chỉ so sánh định vị, phạm vi chức năng và quy ước trách nhiệm được công bố; báo cáo không xếp hạng chất lượng các sản phẩm.

ConferenceSpace là nền tảng thực nghiệm, chưa có bằng chứng về độ trưởng thành trong vận hành, thỏa thuận mức dịch vụ, khả năng hỗ trợ tổ chức hội nghị, quy mô người dùng hoặc mức độ đầy đủ của các chức năng quản lý sự kiện như một sản phẩm thương mại. Các nghiệp vụ cần hoàn thiện thêm gồm cơ chế đăng ký hoặc chọn bài phản biện, kiểm soát hạn chót, một số luồng xử lý bản hoàn thiện sau chấp nhận và các thao tác sau khi có quyết định.

\subsection{Giới hạn của dữ liệu và phương pháp đánh giá}

Các luồng AI chủ yếu được đánh giá trên bài báo tiếng Anh thuộc bộ dữ liệu OpenReview lưu cục bộ. Bộ thử đối sánh phản biện viên sử dụng dữ liệu tổng hợp mô phỏng lược đồ Semantic Scholar, không phải dữ liệu thu thập từ hoạt động phân công thực tế. Hiệu quả của các luồng đối với bài viết tiếng Việt, hội nghị nhỏ, lĩnh vực có cấu trúc bài khác hoặc hội nghị áp dụng chính sách bình duyệt khác chưa được xác nhận. Phương pháp đối sánh chưa được đánh giá bằng dữ liệu phân công thực tế của Chủ tọa; cơ chế phát hiện xung đột lợi ích cũng chưa có tập các cặp có và không có xung đột do chuyên gia gán nhãn.

Một số chức năng AI còn thiếu nhãn chuyên gia, gồm gợi ý chuyên đề, cảnh báo của Submission Gating, mức độ hữu ích của các điểm cần chú ý và chất lượng của bản tổng hợp hỗ trợ quyết định. Các phép chấm TCA và NLI chưa được hiệu chuẩn bằng nhãn chuyên gia; phép chấm có thể bỏ sót mệnh đề đúng được diễn đạt theo cách khác hoặc đánh giá không chính xác các suy luận kỹ thuật phức tạp. Chỉ số Additionality chỉ cho biết mệnh đề được đánh giá là có căn cứ nhưng không trùng trực tiếp với nội dung tham chiếu; chỉ số này không chứng minh thông tin mới hữu ích.

Ba tệp UAT hiện có 91 phản hồi xác nhận đồng ý tham gia và cung cấp đủ dữ liệu để tính các chỉ số chính, gồm 76 phản hồi của Tác giả, 7 phản hồi của Phản biện viên và 8 phản hồi của Chủ tọa. Nhóm tuyển mẫu theo phương pháp thuận tiện; Tác giả chiếm 83,5\% tổng số phản hồi. Do nguồn lực hạn chế và vai trò Chủ tọa hoặc vị trí tương đương đòi hỏi kinh nghiệm chuyên môn, nhóm gặp khó khăn trong việc tuyển người phù hợp và sẵn sàng tham gia. Nhóm chuyển phiếu dành cho Chủ tọa sang hình thức ẩn danh để giảm rào cản tuyển mẫu. Theo quy trình thu thập của nhóm, tám phản hồi này đến từ tám lượt tham gia ẩn danh; tệp xuất ghi cùng một địa chỉ đại diện do lựa chọn ẩn danh, không phải do trùng bản ghi. Tuy nhiên, vì không cấp mã riêng cho từng người tham gia, tệp xuất không cho phép kiểm chứng độc lập tính duy nhất, danh tính hoặc hồ sơ chuyên môn của mẫu Chủ tọa. Vì vậy, kết quả từ nhóm Chủ tọa chỉ cung cấp bằng chứng mô tả ban đầu và không hỗ trợ suy rộng thống kê cho vai trò này. Khảo sát cũng chỉ đo mức độ đồng ý với nhận định rằng AI giúp giảm thao tác hoặc tiết kiệm thời gian; khảo sát không đo thời gian thực hiện nhiệm vụ, gánh nặng nhận thức, chất lượng quyết định hoặc thiên lệch tự động hóa.

\subsection{Giới hạn vận hành, cô lập lỗi và quản trị dữ liệu}

Kiến trúc cung cấp các luồng thao tác thủ công và tách trạng thái nghiệp vụ khỏi kết quả AI, nhưng đề tài chưa thực hiện thử nghiệm chèn lỗi đầu cuối khi yêu cầu đến dịch vụ AI bị timeout hoặc dịch vụ ngừng hoạt động. Do đó, báo cáo chưa chứng minh toàn bộ nghiệp vụ vẫn tiếp tục vận hành và trạng thái không bị thay đổi ngoài ý muốn khi dịch vụ AI gặp sự cố. Một số luồng ghi nhận độ trễ ngoại lệ vượt 100 giây, còn Chatbot Agent có tỷ lệ lỗi gọi công cụ đáng kể. Hệ thống chưa được thử nghiệm khả năng vận hành liên tục hoặc hoạt động phân tán trong thời gian dài.

Cơ chế phân quyền ở cấp ứng dụng giới hạn tài nguyên mà người dùng có thể truy cập, nhưng không thay thế việc kiểm toán toàn bộ vòng đời dữ liệu tại nhà cung cấp AI. Đề tài chưa xác minh chính sách lưu giữ dữ liệu, quyền sử dụng dữ liệu để huấn luyện, khu vực lưu trữ và xử lý dữ liệu hoặc quy trình xử lý sự cố của từng nhà cung cấp. Hệ thống cũng chưa cung cấp cho người dùng nhật ký kiểm toán để theo dõi đầy đủ nguồn dữ liệu, lượt gọi công cụ, đầu ra, nội dung chỉnh sửa và quyết định ghi đè. Vì vậy, chưa có bằng chứng cho thấy ConferenceSpace phù hợp để gửi bản thảo hoặc nội dung phản biện mật đến nhà cung cấp AI bên ngoài. Khi triển khai cho hội nghị thực tế, các luồng này phải bị vô hiệu hóa theo mặc định cho đến khi chính sách hội nghị và thỏa thuận xử lý dữ liệu của nhà cung cấp cho phép rõ ràng~\cite{ch5ref3,ch5ref4}.

\section{Hướng phát triển trong tương lai}

\subsection{Bổ sung bằng chứng còn thiếu}

Ưu tiên đầu tiên là xây dựng tập đánh giá gần với điều kiện vận hành thực tế. Hoạt động đối sánh phản biện viên cần sử dụng dữ liệu phân công lịch sử, nhãn do Chủ tọa cung cấp và tỷ lệ Chủ tọa chấp nhận đề xuất. Cơ chế phát hiện xung đột lợi ích cần tập các cặp có và không có xung đột, bao gồm trường hợp tự phân công, đồng tác giả, cùng đơn vị, quan hệ cố vấn và quan hệ được khai báo, để đo precision, recall và độ phủ theo từng nguồn dữ liệu. Nhiều Phản biện viên và Chủ tọa cần đánh giá các luồng gần bước ra quyết định theo mức độ hữu ích, độ đầy đủ, thời gian đọc, số điểm bất đồng được phát hiện và mức độ ảnh hưởng đến phán đoán.

Các đợt UAT tiếp theo cần mở rộng mẫu ở từng vai trò, đặc biệt là Chủ tọa, thông qua hợp tác với ban tổ chức hội nghị, giảng viên hoặc người từng đảm nhiệm trách nhiệm tương đương. Quy trình nên xác định tiêu chí sàng lọc, cấp mã ẩn danh duy nhất cho từng người tham gia và lưu bằng chứng xác nhận tư cách tách biệt với câu trả lời để vừa bảo vệ danh tính vừa kiểm soát tính hợp lệ của mẫu. Thử nghiệm cần chuẩn hóa nhiệm vụ, đo trực tiếp thời gian thao tác và có điều kiện đối chứng không sử dụng AI. Thiết kế đánh giá cũng cần đo thiên lệch tự động hóa, mức độ người dùng kiểm tra lại nguồn và chất lượng quyết định sau khi nhận gợi ý. Bộ dữ liệu nên được mở rộng sang bài viết tiếng Việt, nhiều lĩnh vực và nhiều chính sách hội nghị.

\subsection{Kiểm chứng khả năng cô lập lỗi và hoàn thiện vận hành}

Hệ thống cần có kịch bản chèn lỗi cho từng dịch vụ AI và Neo4j, gồm yêu cầu bị timeout, dịch vụ của nhà cung cấp gặp lỗi, đầu ra sai cấu trúc và mất kết nối. Thử nghiệm phải xác nhận rằng người dùng vẫn thực hiện được các thao tác thủ công như nộp bài, phân công, gửi bản phản biện và ra quyết định; lỗi AI không tự động thay đổi trạng thái nghiệp vụ; thông báo lỗi nêu rõ cách người dùng cần xử lý.

Các luồng xử lý kéo dài cần sử dụng hàng đợi bất đồng bộ, giới hạn thời gian cho từng giai đoạn, số lần thử lại hữu hạn, cơ chế bảo đảm tác vụ lặp lại không gây sai lệch trạng thái, thông tin tiến độ và khả năng khôi phục. Cơ chế giám sát vận hành cần liên kết yêu cầu, kết quả AI, lượt gọi công cụ, nguồn bằng chứng, phiên bản cấu hình, quyết định xác nhận và thao tác ghi đè. Nhật ký kiểm toán phải hỗ trợ cả việc điều tra lỗi kỹ thuật và xác định trách nhiệm nghiệp vụ.

\subsection{Tăng quyền kiểm soát dữ liệu và hoàn thiện nghiệp vụ}

Đề tài cần đánh giá nhà cung cấp dựa trên điều khoản hợp đồng về dữ liệu, thời gian lưu giữ, quyền sử dụng dữ liệu để huấn luyện, khu vực lưu trữ và xử lý dữ liệu và quy trình xóa dữ liệu. Đối với hội nghị có yêu cầu bảo mật cao, nhóm có thể so sánh phương án triển khai tại chỗ hoặc sử dụng mô hình có trọng số công khai theo các tiêu chí chất lượng, chi phí và công sức vận hành, thay vì mặc định tăng số lượng chức năng AI.

Về sản phẩm, đề tài cần ưu tiên hoàn thiện các nghiệp vụ còn thiếu hoặc chưa đầy đủ, gồm cơ chế đăng ký hoặc chọn bài phản biện, kiểm soát hạn chót, xử lý bản hoàn thiện sau chấp nhận và các thao tác sau khi có quyết định. Nếu ConferenceSpace được phát triển thành hạ tầng học thuật mở, mô hình vận hành nên dựa trên dịch vụ triển khai, hỗ trợ và thỏa thuận mức dịch vụ, không dựa trên việc khai thác dữ liệu bản thảo hoặc nội dung phản biện~\cite{ch5ref8,ch5ref9}.

Tóm lại, đóng góp của ConferenceSpace không nằm ở luận điểm rằng thị trường thiếu AI hoặc việc bổ sung nhiều chức năng AI mặc nhiên tạo ra hệ thống tốt hơn. Trong phạm vi bằng chứng hiện có, đóng góp của đề tài là xây dựng và áp dụng mô hình phân công tác vụ dựa trên hậu quả, khả năng tái tạo và chủ thể chịu trách nhiệm; đồng thời thu thập bằng chứng cho từng tác vụ để xác định những điểm mô hình phù hợp và những điểm hệ thống chưa tuân thủ chính quy ước đã đặt ra. Báo cáo chưa xác nhận tính đúng đắn của toàn bộ vòng đời, mức độ phù hợp chuyên môn của kết quả phân công, chất lượng của cơ chế phát hiện xung đột lợi ích nhiều lớp, mức độ an toàn của dữ liệu tại nhà cung cấp, hiệu lực kiểm soát của con người hoặc khả năng vận hành dài hạn. Các giới hạn này cần được giữ nguyên khi phát triển ConferenceSpace thành hệ thống phục vụ hội nghị thực tế.
